\section{Federated Search-State Coordination}

FL distributes the variables, model weights and evidence that guide NAS across clients and rounds. This theme covers mechanisms that combine, preserve or jointly optimise that state so the search follows a coherent trajectory.

\subsection{Coordinating Distributed Search State}

\subsubsection{Challenge 3.1: Architecture and Model Parameters Form Coupled Global State}

Differentiable NAS jointly optimises model weights and architecture variables. The architecture variables determine which candidate operations contribute to the model, while their gradients depend on the current operation weights. Aggregating only model weights discards local structural decisions, whereas combining architecture variables with unrelated weights separates those decisions from the state that produced them~\cite{fednas_2021,dfnas_2020,feathers_2023}.

\noindent\textbf{Solution 3.1a: Aggregate architecture and model variables together.} Clients can return both variable sets after local search, after which the server aggregates and broadcasts the resulting architecture variables with their corresponding model weights~\cite{fednas_2021,dfnas_2020,feathers_2023}. This maintains one shared differentiable search state across rounds. Parameter-wise aggregation requires clients to use the same supernet topology and assign the same meaning to every architecture variable.

\subsubsection{Challenge 3.2: Locally Searched Architectures Can Be Structurally Incompatible}

Independent local searches can return discrete graphs with different connections, operations or blocks. Their parameters and architecture variables do not occupy a uniform coordinate system, so ordinary parameter or gradient averaging can lose structural information~\cite{multi-granular_fednas_2023,efnas_2024}. Aggregating only the components shared by several local architectures also discards knowledge learned by their non-overlapping components~\cite{f2nas_2024}.

\noindent\textbf{Solution 3.2a: Aggregate recurring subgraphs.} The server can estimate how often connections and operations occur across local graphs and sample a global graph from the resulting distributions~\cite{multi-granular_fednas_2023}. A related strategy directly assembles frequently recurring subgraphs~\cite{efnas_2024}. Both preserve structural information without requiring parameter-wise alignment of complete local architectures.

\noindent\textbf{Solution 3.2b: Transfer knowledge from non-overlapping components.} Parameters of compatible blocks can be aggregated with weights based on their similarity, while predictions from heterogeneous blocks can train the server model through distillation~\cite{f2nas_2024}. This complements the model-agnostic knowledge transfer in Solutions 2.1e and 2.8b. Exchanging compressed predictions similarly enables collaboration between entirely different local topologies~\cite{hetefed_2023}. Distillation avoids direct tensor alignment, but approaches that evaluate predictions on common inputs require a suitable shared dataset.

\subsubsection{Challenge 3.3: Search-Policy Knowledge Can Be Lost between Rounds}

Controller-based reinforcement-learning strategies accumulate search knowledge in a policy that samples candidate architectures. If client state is reset between federated rounds and only candidate weights are retained, the controller loses what earlier evaluations taught it and restarts its local search from a cold state~\cite{efficient_model_training_fednas_2024}.

\noindent\textbf{Solution 3.3a: Persist and redistribute search-policy state.} Clients can upload controller parameters alongside their candidate-weight dictionary. The server retains a selected controller state, aggregates the shared weights and redistributes both for the next round, allowing reinitialised clients to continue from previous search knowledge~\cite{efficient_model_training_fednas_2024}. The cited framework is conceptual and does not empirically compare rules for selecting or combining client controllers.

\subsection{Coordinating Interdependent Search Decisions}

\subsubsection{Challenge 3.4: Architecture Quality Depends on the Federated Training Configuration}

An architecture's apparent quality depends on the configuration under which it is trained and evaluated. Learning rates interact with architectural choices, while the number of local epochs and participating-client fraction change the federated trajectory used to obtain candidate fitness~\cite{feathers_2023,aging_fednas_2025,breastcancer_fednas_2024}. Tuning these variables only after architecture search can therefore invalidate the comparisons that led to the selected architecture.

\noindent\textbf{Solution 3.4a: Coordinate architecture and hyperparameter optimisation.} An alternating procedure can select a training configuration, search an architecture under that configuration and use the result to guide the next hyperparameter decision~\cite{feathers_2023}. Evolutionary search can instead encode the architecture and its training hyperparameters in the same individual and evaluate them together through federated training~\cite{aging_fednas_2025}.

\noindent\textbf{Solution 3.4b: Include FL control variables in the search space.} Black-box strategies can search the number of local epochs and the fraction of active clients alongside architecture and conventional hyperparameter choices. Evaluating each complete configuration through federated training makes candidate fitness reflect the FL procedure with which the architecture will operate~\cite{breastcancer_fednas_2024}.

\subsection{Coordinating Vertically Partitioned Search}

\subsubsection{Challenge 3.6: No Party Can Evaluate an End-to-End Candidate Alone}

In vertical FL, parties hold different features for the same entities and operate separate parts of a joint model. No feature holder can calculate the complete prediction, loss or architecture gradient because its output must be combined with those of the other parties. The local submodels can also have different topologies, making direct parameter aggregation unsuitable~\cite{self-supervised_vertical_fednas_2023}.

\noindent\textbf{Solution 3.6a: Coordinate local searches through intermediate representations.} Each party can retain its own supernet, weights and architecture variables while participating in a split-model forward and backward pass. Feature-holding parties send intermediate outputs to an aggregating party, which computes the joint loss and returns the corresponding gradients. Each party then updates its local weights and architecture variables without requiring compatible local architectures~\cite{self-supervised_vertical_fednas_2023}.

\subsubsection{Challenge 3.7: Coordinated Vertical Search Is Restricted to Aligned Samples}

The joint forward and backward pass requires records that can be matched across all parties. Their aligned intersection may be much smaller than the union of their local datasets, leaving most observations unable to contribute to the supervised architecture-search objective~\cite{self-supervised_vertical_fednas_2023}.

\noindent\textbf{Solution 3.7a: Pretrain local search state on unaligned data.} Before coordinated search, each party can perform self-supervised one-shot NAS on all of its local observations, updating both model weights and architecture variables. These states initialise supervised vertical search on the smaller aligned dataset~\cite{self-supervised_vertical_fednas_2023}. Unaligned data therefore shape the initial representations and structural preferences, while the joint phase supplies the cross-party supervised objective.
